\scanpage{62}
tous nos beaux énoncés devraient être valables, non seulement pour un corps de base $K$ ext. de type fini de $\mathbb{Q}$, mais dès que $K$ est ext. de type fini d'un corps cyclotomique (pas néc. fini sur $\mathbb{Q}$).

Nous pourrions définir les \uline{courbes} de \uline{Poincaré} sur un corps alg. clos $\bar K$ de car $0$, comme étant les courbes isomorphes à des rev. universels de courbes alg. anabéliennes sur $\bar K$ (deux courbes anabéliennes $=\ \bar U,\bar V$ sur $\bar K$ définiss\unsure{e|a}nt des rev. de Poincaré isomorphes, ssi $\exists$ un revêtement fini étale de l'une, isomorphe à un rev. fini étale de l'autre\ldots).

\unsure{E|É}tant donné une courbe de Poincaré\unsure{.} $\unsure{\widetilde{\mathcal U}|\widetilde{\overline{\mathcal U}}}$ sur $\bar K$, \deleted{telle} on définit canoniquement sa complétée\inserted{ $\widehat U$}, de telle façon que $\forall$ autom de $(\widetilde{\mathcal U},\bar K)$ (pas néc. induisant l'identité sur $\bar K$) induise un autom. de $\widehat{\mathcal U}$.
Ceci posé :
\margin{61}
